\documentclass[11pt,a4paper]{article}
\usepackage[T1]{fontenc}
\usepackage[utf8]{inputenc}
\usepackage[brazil]{babel}
\usepackage[margin=1.8cm]{geometry}
\usepackage{amsmath,amssymb}
\usepackage{booktabs}
\usepackage{enumitem}
\usepackage{caption}
\usepackage{microtype}
\usepackage[compact]{titlesec}
\usepackage[hidelinks]{hyperref}

\titlespacing*{\section}{0pt}{0.9em}{0.4em}
\captionsetup{font=small,labelfont=bf}
\setlist{nosep,leftmargin=1.4em}
\setlength{\parskip}{0.25em}
\pagestyle{plain}

\newcommand{\cmax}{C_{\max}}

\begin{document}

\noindent{\LARGE\bfseries Otimização do Sequenciamento na Produção de Baterias:\\[2pt] Abordagem via RCPSP}\\[6pt]
{\bfseries Etapa 6: Resultados do Método Aproximado (Metaheurística)}\\
Cristian Madeira (Pós), Otávio Novais (162962), Thaís Melo (164934)

\section*{Configuração dos experimentos}

O método aproximado é o RKO (\emph{Random-Key Optimizer}) \cite{rko}, na implementação C++ dos
autores, sem alterações, com um decodificador próprio para o RCPSP: o vetor de chaves gera uma
lista de prioridade topologicamente válida (a cada passo, a atividade elegível de menor chave),
convertida em cronograma pelo SSGS e melhorada pelo \emph{forward-backward improvement} (FBI)
\cite{valls}. O RKO executa nove metaheurísticas em paralelo (BRKGA, SA, GRASP, ILS, VNS, PSO, GA,
LNS e BRKGA-CS), que compartilham um conjunto de soluções de elite, com os parâmetros padrão do
framework. Foram usados dois critérios de parada:

\begin{itemize}
  \item \textbf{Tempo:} 10~s por execução no J30 e 30~s no J60.
  \item \textbf{Orçamento de 50.000 cronogramas} (decodificações, somadas entre as threads), o
        critério usual da literatura \cite{kolisch}. A criação do conjunto inicial de soluções
        consome cerca de 10.800 cronogramas no J30 e 21.600 no J60, igual para todos os métodos.
\end{itemize}

Em ambos, a execução também termina ao atingir o ótimo conhecido. Cada instância foi resolvida
10 vezes, no mesmo notebook da Etapa~5 (Intel Core i5-1135G7, 16~GB). Às dez instâncias J30 da
Etapa~5 somaram-se cinco J60 com $RS = 0{,}2$ e ótimo conhecido, também resolvidas pelo Gurobi
com 600~s. A qualidade é medida pelo desvio relativo ao ótimo, RPD $= 100\,(\cmax - C^*)/C^*$.

\section*{Exato $\times$ aproximado}

\begin{table}[ht]
\centering
\small
\setlength{\tabcolsep}{4.3pt}
\caption{Gurobi (Etapa 5) e RKO nas mesmas instâncias. RKO: melhor e médio $\cmax$ em 10 execuções e tempo médio por execução.}
\label{tab:exato-rko}
\begin{tabular}{lrrrrrrrrr}
\toprule
 & & \multicolumn{2}{c}{Gurobi (600 s)} & \multicolumn{3}{c}{RKO, tempo (10 s / 30 s)} & \multicolumn{3}{c}{RKO, 50.000 cronogramas} \\
\cmidrule(lr){3-4}\cmidrule(lr){5-7}\cmidrule(lr){8-10}
Instância & $C^*$ & $\cmax$ & Tempo (s) & Melhor & Médio & Tempo (s) & Melhor & Médio & Tempo (s) \\
\midrule
% linhas geradas por relatorios/etapa6/gera_tabela.py
j301\_1 & 43 & \textbf{43} & 0,1 & \textbf{43} & 43,0 & 0,10 & \textbf{43} & 43,0 & 0,09 \\
j309\_2 & 92 & \textbf{92} & 286,5 & \textbf{92} & 92,0 & 0,12 & \textbf{92} & 92,0 & 0,11 \\
j3013\_1 & 58 & 59$^\dagger$ & 600,1 & \textbf{58} & 58,3 & 5,70 & 60 & 60,0 & 0,28 \\
j3025\_4 & 81 & \textbf{81} & 242,8 & \textbf{81} & 81,8 & 8,98 & 82 & 82,0 & 0,26 \\
j3029\_8 & 80 & 81$^\dagger$ & 600,1 & 81 & 81,0 & 10,20 & 82 & 82,0 & 0,33 \\
j3037\_7 & 92 & \textbf{92} & 40,9 & \textbf{92} & 92,0 & 0,11 & \textbf{92} & 92,0 & 0,10 \\
j3041\_10 & 99 & \textbf{99} & 165,7 & \textbf{99} & 99,0 & 0,14 & \textbf{99} & 99,0 & 0,12 \\
j3045\_6 & 129 & \textbf{129}$^\dagger$ & 600,1 & \textbf{129} & 129,0 & 0,15 & \textbf{129} & 129,0 & 0,13 \\
j3014\_2 & 53 & \textbf{53} & 51,0 & \textbf{53} & 53,0 & 0,12 & \textbf{53} & 53,0 & 0,11 \\
j3030\_2 & 68 & \textbf{68} & 6,0 & \textbf{68} & 68,0 & 0,57 & \textbf{68} & 68,4 & 0,21 \\
\midrule
j605\_4 & 72 & 77$^\dagger$ & 600,1 & \textbf{72} & 72,5 & 22,43 & 73 & 74,4 & 0,84 \\
j609\_4 & 87 & 92$^\dagger$ & 600,1 & 88 & 88,5 & 30,80 & 90 & 91,0 & 0,94 \\
j6021\_4 & 95 & 96$^\dagger$ & 600,1 & \textbf{95} & 95,0 & 3,71 & 98 & 99,0 & 0,74 \\
j6037\_2 & 95 & 96$^\dagger$ & 600,1 & 96 & 96,0 & 30,60 & 97 & 97,0 & 0,69 \\
j6041\_7 & 132 & 147$^\dagger$ & 600,1 & 134 & 134,0 & 30,71 & 137 & 138,5 & 0,84 \\
\midrule
RPD / tempo J30 & & 0,30\% & 259,3 & 0,12\% & 0,28\% & 2,62 & 0,72\% & 0,78\% & 0,17 \\
RPD / tempo J60 & & 5,23\% & 600,1 & 0,74\% & 1,00\% & 23,65 & 2,78\% & 3,83\% & 0,81 \\
\bottomrule
\end{tabular}
\par\smallskip
\begin{minipage}{0.97\textwidth}\footnotesize
Em negrito: igual ao ótimo $C^*$. $^\dagger$: limite de tempo atingido, sem prova de otimalidade.
Últimas linhas: RPD médio de cada coluna de $\cmax$ e tempo médio, por conjunto.
\end{minipage}
\end{table}

\textbf{J30.} Os dois métodos chegam à mesma qualidade, mas o RKO é cerca de 100 vezes mais
rápido: RPD médio de 0,28\% em 2,6~s por execução, contra 0,30\% do Gurobi em 259~s. Em j3013\_1,
o RKO atingiu o ótimo (58), enquanto o Gurobi terminou em 59. Em j3029\_8, nenhum dos dois
chegou ao ótimo (80). A vantagem do método exato é a prova de otimalidade, obtida em 7 das 10
instâncias.

\textbf{J60.} O Gurobi não provou a otimalidade de nenhuma instância em 600~s, com RPD médio de
5,2\% (até 11,4\% em j6041\_7: 147 contra o ótimo 132). Com 30~s, o RKO obteve RPD médio de 1,0\%,
igual ou melhor que o Gurobi nas cinco instâncias e no ótimo em j605\_4 e j6021\_4. Mesmo limitado
a 50.000 cronogramas (menos de 1~s), o RKO supera o Gurobi em 3 das 5 instâncias, com RPD médio de
3,8\%. O ganho do método aproximado cresce com o tamanho, como esperado para a formulação
tempo-indexada, cujo número de variáveis acompanha o produto entre atividades e horizonte.

\section*{Comparação entre estratégias de busca}

Para avaliar a observação de que algoritmos genéticos têm dificuldade neste problema, o GA e o
BRKGA \cite{brkga} do mesmo framework foram executados isoladamente, com o mesmo decodificador,
com e sem o FBI.

\begin{table}[ht]
\centering
\small
\caption{RPD médio (\%) nas 15 instâncias, 10 execuções cada; $p$: teste de Friedman entre os três métodos.}
\label{tab:busca}
\begin{tabular}{llrrrr}
\toprule
Critério de parada & Decodificador & GA & BRKGA & RKO & $p$ \\
\midrule
Tempo (10 s / 30 s) & SSGS + FBI & 0,50 & 0,49 & 0,52 & 0,81 \\
                    & Só SSGS    & 0,95 & 1,03 & 0,82 & 0,12 \\
50.000 cronogramas  & SSGS + FBI & 1,71 & 1,81 & 1,80 & 0,27 \\
                    & Só SSGS    & 4,26 & 4,84 & 4,43 & 0,04 \\
\bottomrule
\end{tabular}
\end{table}

Com o FBI, os três métodos são estatisticamente equivalentes nos dois critérios. A única diferença
significativa aparece sem o FBI e com orçamento fixo: o GA supera o BRKGA (Wilcoxon com correção
de Holm, $p = 0{,}02$), e o RKO fica entre os dois. A dificuldade esperada do GA, portanto, não se
confirmou quando ele usa a mesma representação e o mesmo decodificador. Uma explicação provável
para o desempenho do BRKGA é o orçamento: o framework divide os cronogramas entre nove cópias
paralelas, e, com a população padrão de 987 indivíduos, cada cópia do BRKGA faz só 3 a 4 gerações,
contra 8 a 11 do GA (população de 400). Os parâmetros não foram calibrados para esse orçamento, o
que limita a conclusão. O fator de maior efeito é o decodificador: o FBI reduz o RPD médio de
1,8--2,1\% para 0,7--0,8\% no J30 e de 9,1--10,4\% para 3,6--4,1\% no J60 com orçamento fixo
($p < 0{,}001$ nos três métodos).

\textbf{Síntese.} Para instâncias pequenas, o modelo exato entrega o ótimo com prova, a um custo de
minutos. Para instâncias maiores, o RKO encontra soluções melhores em segundos. O RKO será a
metaheurística do artigo, com o GA e o BRKGA como referência de comparação.

\vspace{-0.6em}
{\scriptsize
\renewcommand{\refname}{\normalsize Referências}
\begin{thebibliography}{9}\setlength{\itemsep}{0pt}
\bibitem{rko} Chaves, A. A. et al. A random-key optimizer for combinatorial optimization. \emph{Journal of Heuristics}, 31(4), 2025.
\bibitem{brkga} Gonçalves, J. F.; Resende, M. G. C. Biased random-key genetic algorithms for combinatorial optimization. \emph{Journal of Heuristics}, 17(5):487--525, 2011.
\bibitem{kolisch} Kolisch, R.; Hartmann, S. Experimental investigation of heuristics for resource-constrained project scheduling: an update. \emph{European Journal of Operational Research}, 174(1):23--37, 2006.
\bibitem{valls} Valls, V.; Ballestín, F.; Quintanilla, S. Justification and RCPSP: a technique that pays. \emph{European Journal of Operational Research}, 165(2):375--386, 2005.
\end{thebibliography}}

\end{document}
